\documentclass[a4paper,11pt]{article}
\usepackage{fancyhdr,graphics,graphicx,xy}
\usepackage{blkarray}
\usepackage{ifthen}
\usepackage{amssymb,amsmath,mathtools,color}
\pagestyle{fancy}
\usepackage{xparse}
\usepackage{nicematrix}


%\textwidth = 6.5 in
%\textheight = 9 in
%\oddsidemargin = 0.0 in
%\evensidemargin = 0.0 in
%\topmargin = 0.0 in
%\headheight = 0.0 in
%\headsep = 0.0 in
%\parskip = 0.2in
%\parindent = 0.0in


%%%%%%%%%%%%%%%%%%%%%%%%%%%%%%%%%%%%%%%%%%
%%% General math macros
%%%%%%%%%%%%%%%%%%%%%%%%%%%%%%%%%%%%%%%%%%

\newcommand{\tf}{\tfrac} 
\renewcommand{\Re}{\mathrm{Re}}

\newcommand{\Def}{\coloneqq}
\newcommand{\lgg}{\operatorname{lg}}
\newcommand{\R}{\mathbb{R}}
\newcommand{\C}{\mathbb{C}}
\newcommand{\D}{\mathbb{D}}
\newcommand{\A}{\mathcal{A}}
\newcommand{\F}{\mathcal{F}}
\newcommand{\Q}{\mathbb{Q}}
\newcommand{\N}{\mathbb{N}} 
\newcommand{\Quat}{\mathbb{J}}
\newcommand{\Z}{\mathbb{Z}}
\newcommand{\T}{\mathbb{T}}
\newcommand{\U}{\mathcal{U}}
\newcommand{\LP}{\operatorname{L}}
\newcommand{\abs}[1]{\left| {#1} \right|}
\newcommand{\sgn}{\operatorname{sgn}}

\newcommand{\tendsto}{\to}
\newcommand{\tendto}{\to}

%%%%%%%%%%%%%%%%%%%%%%%%%%%%%%%%%%%%%%%%%%
%%% Probability Macros
%%%%%%%%%%%%%%%%%%%%%%%%%%%%%%%%%%%%%%%%%%
\newcommand{\stdcg}{N_\bbc(0,1)}

\newcommand{\expec}{\mathbb{E}}
\newcommand{\expect}{\mathbb{E}}
\newcommand{\Exp}{\mathbb{E}}
\newcommand{\E}{\mathbb{E}}
\newcommand{\prob}{\mathbb{P}}
\newcommand{\filt}{\mathscr{F}}
\newcommand{\Prob}{\prob}
\renewcommand{\Pr}{\prob}
\newcommand{\Ent}{\operatorname{Ent}}
\newcommand{\Var}{\operatorname{Var}}
\newcommand{\Cov}{\operatorname{Cov}}
\newcommand{\dtv}{d_{TV}}
\newcommand{\weakto}{\Rightarrow}
\newcommand{\lawof}{\mathscr{L}}
\newcommand{\Bernoulli}{\operatorname{Bernoulli}}
\newcommand{\Binomial}{\operatorname{Binom}}
\newcommand{\Poisson}{\operatorname{Poisson}}
\newcommand{\Normal}{\operatorname{Normal}}
\newcommand{\Unif}{\operatorname{Unif}}
\newcommand{\As}{\ensuremath{\text{a.s.}}}
\newcommand{\distequals}{\overset{\mathrm{d}}{=}}


\newcommand{\asstnumber}{5}
\newcommand{\assigntop}{
	\setlength{\headheight}{33.40004pt}
  \fancyhead[L]{}
  \fancyhead[C]{\Large Math 356 -- Fall 2023\\ Assignment \asstnumber}
  \fancyhead[R]{ }
}

\newtheorem{theorem}{Theorem}
\newtheorem{corollary}[theorem]{Corollary}
\newtheorem{definition}{Definition}
\newcommand{\p}[1]{\mathbb{P}(#1)}
\newcommand{\px}[1]{\mathbb{P}_x(#1)}
\newcommand{\py}[1]{\mathbb{P}_y(#1)}
\newcommand{\I}[1]{\mathbf{1}_{#1}}
\newcommand{\e}{\mathbb{E}}

\begin{document}
\assigntop

\section{ Question 1}
Let $(X_1,X_2,X_3,X_4) \sim Mult(n,4,p_1,p_2,p_3,p_4)$. We are able to find the distribution of $X_1+X_2$, be realizing that $X_1$ counts the number of times that $1$ occured and $X_2$ counts the number of times that $2$ occured. The probability that $1$ occurs is $p_1$ and the probability that $2$ occurs is $p_2$, since these are independent trials (by def of the multinomial distribution), the probability that $1$ \textbf{or} $2$ occurs is $p_1 + p_2$ therefore $X_1 + X_2 \sim Bin(n,p_1+p_2)$.

\clearpage

\section{ Question 2}
Suppose $X$ and $Y$ are jointly continuous with joint probability density function 
$$f(x,y) = ce^{\frac{x^2}{2} - \frac{(x-y)^2}{2}}, x,y \in (-\infty, +\infty)$$
	for some constant $c$
	\subsection*{(a)}
	For this to be a joint probability density function, it's integral over the reals must be $1$. Therefore, we can find the constant $c$ by computing the integral 
	\begin{align*}
		\int_{-\infty}^\infty e^{-x^2/2 - (x-y)^2/2} dy &= e^{-x^2/2} \int_{-\infty}^{\infty} e^{-\frac{(x-y)^2}{2}}dy \\
								&= \sqrt{2\pi} e^{-x^2/2} \underbrace{\int_{-\infty}^{\infty} \frac{1}{\sqrt{2\pi}} e^{\frac{(x-y)^2}{2}}}_{\text{$N(x,1)$ integrates to $1$}} dy &= \sqrt{2\pi}e^{\frac{-x^2}{2}}
	\end{align*}
	Completing the integral over $dx$, we get 
	\begin{align*}
		\underbrace{\int_{0}^{+\infty}}_{\text{ $e$ is non negative}} \sqrt{2\pi}e^{\frac{-x^2}{2}}dx  &= 2\pi \int_{0}^\infty \frac{1}{\sqrt{2\pi}} e^{-x^2}{2}dx = 2\pi
	\end{align*}
therefore, for the integral over the entire reals to be $1$, we need to choose $c = \frac{1}{2\pi}$
	\subsection*{(b)}
	Computing $f_X$ was already done in part (a), by integrating over $dy$ first, we get 
	$$
		f_X = \frac{1}{\sqrt{2\pi}} e^{\frac{-x^2}{2}}
	$$
	We want to compute $f_Y$ by integrating over $dx$ only, 
	\begin{align*}
		f_Y(y) = \frac{1}{\sqrt{2\pi}}\int_{-\infty}^{+\infty} \frac{1}{\sqrt{2\pi}} e^{\frac{-x^2}{2} - \frac{-(x-y)^2}{2}}
	\end{align*}
	We want to get this integral in a form that we recognize, to do this observe the following trick
	\begin{align*}
		\frac{-x^2}{2} - \frac{(x-y)^2}{2} &= \frac{-x^2}{2} - \frac{x^2-2xy+y^2}{2}\\
						   &= -(x^2 - xy + \frac{y^2}{2})\\
						   &= -( x^2 - xy + \frac{y^2}{4} + \frac{y^2}{4} )\\
						   &= -(x^2 - y/2)^2 - y^2/4
	\end{align*}
	then the integral becomes
	\begin{align*}
		\frac{1}{\sqrt{2\pi}} \int_{-\infty}^{\infty} \frac{1}{\sqrt{2\pi}} e^{-(x^2-y/2)^2 - y^2/4} dx &= \frac{1}{\sqrt{4\pi}}e^{-y^2/4} \int \frac{1}{\sqrt{2\pi \cdot \sqrt{1/2}^2}} e^{\frac{-(x^2-y/2)^2}{2 \cdot \sqrt{1/2}^2}}\\
		&= \frac{1}{\sqrt{4\pi}} e^{-y^2/4}
	\end{align*}
	Where we can recognize that the before last line contains the integral of the PDF of a normal distribution with mean $y/2$ and standard deviation $\sqrt{1/2}$. And we can recognize that $f_Y$ is normal distributed with mean $0$ and variance $2$. 
	\subsection*{(c)} It's clear that $X$ and $Y$ are not independent, since $f(x,y) \neq f_X \cdot f_Y$
	

	


\clearpage
\section*{Question 3}
$X$ and $Y$ are independent, using the change of variables $u = g(x,y) = \frac{x}{x+y}$ and $v = h(x,y) = x+y$, then we can find the inverse change of variable functions to be $p(u,v) = uv$ and $q(u,v) = 	\frac{x+y-x}{x+y} \cdot x+y = (1-u)v$
where 
$$ \det \text{Jacobian}(u,v) = \det \begin{vmatrix}
	v & u \\
	-v & 1-u
\end{vmatrix} = v$$ 
Then we have the formula as presented in class, where $(U,V)$ has joint density as given by 
\begin{align*}
	f_{U,V} &= f(uv, v(1-u))v = f(uv)f(v(1-u)) v \\
		&= \frac{\lambda^r (uv)^{r-1}}{\Gamma(r)} e^{-\lambda u v} \frac{\lambda^s(v(1-u))^{s-1}}{\Gamma(s)} e^{-\lambda(v(1-u))}v\\
		&= \frac{\Gamma(r+s)}{\Gamma(r)\Gamma(s)} u^{r-1}(1-u)^{s-1} \cdot \frac{1}{\Gamma(r+s)} \lambda^{r+s}v^{(r+s)-1}e^{-\lambda v}
\end{align*}
This is exactly the product of the density function of $B \sim Beta(r,s)$ and $G \sim Gamma(r+s,\lambda)$ for subquestion $(b)$.
\subsection*{(b)} 
Since this is exactly the inverse transformation of part $(a)$, then $(X,Y)$'s joint distribution is exactly the product of the density function of $X$ and $Y$ in part (a). To be pedantic, denote $x = b \cdot g$ and $y = (1-b)g$, then the inverse functions are 
$b = \frac{x}{x+y}$ and $g = x+y$, where the determinant of the Jacobian is 
$$ \det \text{Jacobian}(x,y) = \det \begin{vmatrix}
	\frac{y}{(x+y)^2} & - \frac{x}{(x+y)^2} \\ 
	1 & 1
\end{vmatrix} =  \frac{x+y}{(x+y)^2} = \frac{1}{x+y}$$
\clearpage
then we can compute the joint density as 
\begin{align*}
	f_{X,Y} (x,y) &= \frac{\Gamma(r+s)}{\Gamma(r)\Gamma(s)} (\frac{x}{x+y})^{r-1} (1 - \frac{x}{x+y})^{s-1} \cdot \frac{1}{\Gamma(r+s)}\lambda^{r+s} (x+y)^{(r+s)-1}e^{-\lambda(x+y)} \frac{1}{x+y} \\
		      &= \frac{x^{r-1}\lambda^r}{\Gamma(r)} e^{-\lambda x} \frac{y^{s-1}\lambda^s}{\Gamma(s)} e^{-\lambda y}
	\end{align*}
	Where we recognize it to be the product of the density functions of the two independent random variables $X \sim Gamma(r , \lambda)$ and $Y \sim Gamma(s,\lambda)$
\section*{Question 5}
First of all, since $X_1,X_2,X_3$ are independently $Exp(\lambda)$ distributed random variables and they all have the same probability distribution, by theorem 7.20 and definition 7.19 in the textbook, then they are exchangeable. Since they are jointly continuous, then $\Prob(X_i = X_j) = 0$ for any $i \neq j$. This follows from the fact that the space $D = \{(x_1,x_2,x_3) \colon x_i = x_j \text{ such that $i \neq j$} \}$, when intergrated over, that means $\Prob((X_1,X_2,X_3) \in D)$ has value $0$.  Then taking the complement of that $1 - \Prob(X_i = X_j \colon i \neq j) = \Prob(X_1 \neq X_2 \neq X_3) = 1$. To solve this question, we are interested in the events where $X_1 < X_2 < X_3$, $X_1 < X_3 < X_2$, $X_2 < X_1 < X_3$, $X_2 < X_3 < X_1$, $X_3 < X_1 < X_2$ and $X_3 < X_2 < X_1$, where we list all $6$ permutation of $X_1, X_2$ and $X_3$ in these inequalities. The probability of each of those, by the exchangeability of $X_1, X_2$ and $X_3$ are equal. Moreover, they sum up to $\Prob(X_1 \neq X_2 \neq X_3) = 1$, therefore, $\Prob(X_1 < X_2 < X_3) = \frac{1}{6}$. 
\end{document}

